% GENERATED FILE. Edit canonical lessons, then run npm run generate:books.
% canonical-source-hash: fnv1a64:e51f2ea3dad17fa4
% canonical-lessons: RU-C77-city, RU-C77-street, RU-C77-shop, RU-C77-market, RU-C77-school

\chapter{City, Street, Shop, Market, School}
\label{ch:ru-city-street-shop-market-school}
\hlchaptermodality{77}

\begin{chapteropening}
\textbf{By the end of this chapter:} \emph{I can name a city, a street, a shop, a market and a school.}

Complete the last lesson of chapter 77: I can name a city, a street, a shop, a market and a school.
\end{chapteropening}

\section[\texorpdfstring{\ru{город} (górod)}{город (górod)}]{\ru{город} (górod) — a city, a town}

\label{lesson:RU-C77-city}



\begin{quote}
Before the new one: say the Russian for a thousand, then the Russian for a country.
\end{quote}

\subsection*{You'll want to know: \ru{город}}
\textbf{\ru{город}} (\emph{górod}) — \textquotedblleft{}a city, a town\textquotedblright{}.

\begin{cousinweb}
\textbf{\ru{город}} $\leftarrow$ Proto-Slavic \emph{*gordŭ}, \textquotedblleft{}an enclosure, a fort\textquotedblright{} — inherited. Vasmer lists Lithuanian \emph{gardas}, \textquotedblleft{}a fold\textquotedblright{}, Gothic \emph{gards}, \textquotedblleft{}house\textquotedblright{}, and Sanskrit \emph{gṛhá-}, \textquotedblleft{}house\textquotedblright{}: English \emph{yard} is a cousin. Church Slavonic has \emph{\ru{град}}, \textquotedblleft{}city\textquotedblright{}, as in \ru{Петроград}, \textquotedblleft{}Peter's city\textquotedblright{}.
\end{cousinweb}

\begin{grammarlens}[title={What you've built}]
One more word for this chapter.
\end{grammarlens}

\subsection*{Your turn}
\begin{itemize}
\raggedright
  \item \emph{Say it:} \emph{górod}
  \item \emph{Say it:} \emph{górod}, once more
  \item \emph{Say it:} say \emph{górod}
  \item \emph{Recall:} say the Russian for a thousand, then the Russian for a country, then say \emph{górod} again
\end{itemize}

\begin{tcolorbox}[breakable,colback=teal!4,colframe=teal!35!black,title={Before you move on}]
What does \ru{город} mean? (\textquotedblleft{}A city, a town\textquotedblright{}.) Say it once more.
\end{tcolorbox}
\section[\texorpdfstring{\ru{улица} (úlitsa)}{улица (úlitsa)}]{\ru{улица} (úlitsa) — a street}

\label{lesson:RU-C77-street}



\begin{quote}
Before the new one: say the Russian for a country, then the Russian for a city.
\end{quote}

\subsection*{You'll want to know: \ru{улица}}
\textbf{\ru{улица}} (\emph{úlitsa}) — \textquotedblleft{}a street\textquotedblright{}.

\begin{cousinweb}
\textbf{\ru{улица}} $\leftarrow$ Proto-Slavic \emph{*ulica}, built on \emph{*ula}, \textquotedblleft{}a hollow, a channel\textquotedblright{} — inherited. Vasmer relates it to \emph{\ru{улей}}, \textquotedblleft{}beehive\textquotedblright{}, and Greek \emph{aulós}, \textquotedblleft{}pipe\textquotedblright{}, \emph{aul\'{\={o}}n}, \textquotedblleft{}glen\textquotedblright{}: a street as a passage between houses.
\end{cousinweb}

\begin{grammarlens}[title={What you've built}]
One more word for this chapter.
\end{grammarlens}

\subsection*{Your turn}
\begin{itemize}
\raggedright
  \item \emph{Say it:} \emph{úlitsa}
  \item \emph{Say it:} \emph{úlitsa}, once more
  \item \emph{Say it:} say \emph{úlitsa}
  \item \emph{Recall:} say the Russian for a country, then the Russian for a city, then say \emph{úlitsa} again
\end{itemize}

\begin{tcolorbox}[breakable,colback=teal!4,colframe=teal!35!black,title={Before you move on}]
What does \ru{улица} mean? (\textquotedblleft{}A street\textquotedblright{}.) Say it once more.
\end{tcolorbox}
\section[\texorpdfstring{\ru{магазин} (magazín)}{магазин (magazín)}]{\ru{магазин} (magazín) — a shop}

\label{lesson:RU-C77-shop}



\begin{quote}
Before the new one: say the Russian for a city, then the Russian for a street.
\end{quote}

\subsection*{You'll want to know: \ru{магазин}}
\textbf{\ru{магазин}} (\emph{magazín}) — \textquotedblleft{}a shop\textquotedblright{}.

\begin{cousinweb}
\textbf{\ru{магазин}} — first in Peter the Great's time; it came through German \emph{Magasin} from French \emph{magasin}, \textquotedblleft{}a storehouse\textquotedblright{} (Vasmer). The French word goes back, through Italian, to Arabic \emph{makhāzin}, \textquotedblleft{}storehouses\textquotedblright{}.
\end{cousinweb}

\begin{grammarlens}[title={What you've built}]
One more word for this chapter.
\end{grammarlens}

\subsection*{Your turn}
\begin{itemize}
\raggedright
  \item \emph{Say it:} \emph{magazín}
  \item \emph{Say it:} \emph{magazín}, once more
  \item \emph{Say it:} say \emph{magazín}
  \item \emph{Recall:} say the Russian for a city, then the Russian for a street, then say \emph{magazín} again
\end{itemize}

\begin{tcolorbox}[breakable,colback=teal!4,colframe=teal!35!black,title={Before you move on}]
What does \ru{магазин} mean? (\textquotedblleft{}A shop\textquotedblright{}.) Say it once more.
\end{tcolorbox}
\section[\texorpdfstring{\ru{рынок} (rýnok)}{рынок (rýnok)}]{\ru{рынок} (rýnok) — a market}

\label{lesson:RU-C77-market}



\begin{quote}
Before the new one: say the Russian for a street, then the Russian for a shop.
\end{quote}

\subsection*{You'll want to know: \ru{рынок}}
\textbf{\ru{рынок}} (\emph{rýnok}) — \textquotedblleft{}a market\textquotedblright{}.

\begin{cousinweb}
\textbf{\ru{рынок}} came through Polish \emph{rynek} and Czech \emph{rynk}, \textquotedblleft{}town square\textquotedblright{}, from Middle High German \emph{rinc}, \textquotedblleft{}a ring, a square\textquotedblright{} (Vasmer): the market sat in the ring of the square. German \emph{Ring} is the same word.
\end{cousinweb}

\begin{grammarlens}[title={What you've built}]
One more word for this chapter.
\end{grammarlens}

\subsection*{Your turn}
\begin{itemize}
\raggedright
  \item \emph{Say it:} \emph{rýnok}
  \item \emph{Say it:} \emph{rýnok}, once more
  \item \emph{Say it:} say \emph{rýnok}
  \item \emph{Recall:} say the Russian for a street, then the Russian for a shop, then say \emph{rýnok} again
\end{itemize}

\begin{tcolorbox}[breakable,colback=teal!4,colframe=teal!35!black,title={Before you move on}]
What does \ru{рынок} mean? (\textquotedblleft{}A market\textquotedblright{}.) Say it once more.
\end{tcolorbox}
\section[\texorpdfstring{\ru{школа} (shkóla)}{школа (shkóla)}]{\ru{школа} (shkóla) — a school}

\label{lesson:RU-C77-school}



\begin{quote}
Before the new one: say the Russian for a shop, then the Russian for a market.
\end{quote}

\subsection*{You'll want to know: \ru{школа}}
\textbf{\ru{школа}} (\emph{shkóla}) — \textquotedblleft{}a school\textquotedblright{}. Where the \textbf{\ru{учитель}} and the \textbf{\ru{ученик}} meet.

\begin{cousinweb}
\textbf{\ru{школа}} — recorded from 1388; it came through Polish \emph{szkoła} from Latin \emph{schola}, from Greek \emph{schol\'{\={e}}}, \textquotedblleft{}leisure\textquotedblright{}, then \textquotedblleft{}lecture, school\textquotedblright{} (Vasmer).
\end{cousinweb}

\begin{grammarlens}[title={What you've built}]
One more word for this chapter.
\end{grammarlens}

\subsection*{Your turn}
\begin{itemize}
\raggedright
  \item \emph{Say it:} \emph{shkóla}
  \item \emph{Say it:} \emph{shkóla}, once more
  \item \emph{Say it:} say \emph{shkóla}
  \item \emph{Recall:} say the Russian for a shop, then the Russian for a market, then say \emph{shkóla} again
\end{itemize}

\begin{tcolorbox}[breakable,colback=teal!4,colframe=teal!35!black,title={Before you move on}]
What does \ru{школа} mean? (\textquotedblleft{}A school\textquotedblright{}.) Say it once more.
\end{tcolorbox}
